% Bagian: sets-functions-relations
% Bab: functions
% Subbagian: composition

\documentclass[../../../include/open-logic-section]{subfiles}

\begin{document}

\olfileid[id]{sfr}{fun}{cmp}
\olsection{Komposisi Fungsi}

\begin{explain}
\oliflabeldef{sfr:fun:inv:sec}{Dalam \olref[inv]{sec}, kita melihat bahwa
invers~$f^{-1}$ dari !!a{bijection}~$f$ juga merupakan fungsi. Operasi lain
pada fungsi adalah komposisi: k}{K}ita dapat mendefinisikan fungsi baru dengan
mengomposisikan dua fungsi, $f$ dan~$g$, yaitu dengan terlebih dahulu
menerapkan $f$, kemudian~$g$. Tentu saja, hal ini hanya mungkin jika daerah
hasil dan domainnya bersesuaian, yaitu daerah hasil~$f$ harus merupakan
himpunan bagian dari domain~$g$. \oliflabeldef{sfr:rel:ops:sec}{Operasi pada
fungsi ini merupakan padanan operasi hasil kali relatif pada relasi dari
\olref[rel][ops]{sec}.}{}

Sebuah diagram dapat membantu menjelaskan gagasan komposisi. Dalam
\olref{fig:composition}, kita menggambarkan dua fungsi $f \colon A \to B$ dan
$g \colon B \to C$ beserta komposisinya~$(\comp{f}{g})$. Fungsi
$(\comp{f}{g}) \colon A \to C$ memasangkan setiap !!{element} dari~$A$ dengan
!!a{element} dari~$C$. Kita menentukan !!{element} mana dari~$C$ yang
dipasangkan dengan !!a{element} dari $A$ sebagai berikut: jika diberikan
masukan $x \in A$, mula-mula terapkan fungsi $f$ pada~$x$, yang menghasilkan
suatu $f(x) = y \in B$, kemudian terapkan fungsi $g$ pada~$y$, yang
menghasilkan suatu $g(f(x)) = g(y) = z \in C$.
\begin{figure}
  \olasset[2\olphotowidth]{assets/diagrams/composition.tikz}
  \caption{Komposisi $g \circ f$ dari dua fungsi $f$ dan~$g$.}
  \ollabel{fig:composition}
\end{figure}
\end{explain}

\begin{defn}[Komposisi]
Misalkan $f\colon A \to B$ dan $g\colon B \to C$ fungsi. \emph{Komposisi}
dari $f$ dengan~$g$ adalah $\comp{f}{g} \colon A \to C$, dengan
$(\comp{f}{g})(x) = g(f(x))$.
\end{defn}

\begin{ex}
Tinjau fungsi $f(x) = x + 1$ dan $g(x) = 2x$. Karena
$(\comp{f}{g})(x) = g(f(x))$, untuk setiap masukan~$x$ Anda harus terlebih
dahulu mengambil penerusnya, kemudian mengalikan hasilnya dengan dua. Jadi,
komposisi keduanya diberikan oleh $(\comp{f}{g})(x) = 2(x+1)$.
\end{ex}

\begin{prob}
Tunjukkan bahwa jika $f \colon A \to B$ dan $g \colon B \to C$ keduanya
!!{injective}, maka $\comp{f}{g}\colon A \to C$ bersifat !!{injective}.
\end{prob}

\begin{prob}
Tunjukkan bahwa jika $f \colon A \to B$ dan $g \colon B \to C$ keduanya
!!{surjective}, maka $\comp{f}{g}\colon A \to C$ bersifat !!{surjective}.
\end{prob}

\begin{prob}
Andaikan $f \colon A \to B$ dan $g \colon B \to C$. Tunjukkan bahwa graf
dari $\comp{f}{g}$ adalah $R_f \mid R_g$.
\end{prob}

\end{document}
